% GENERATED FILE. Edit canonical lessons, then run npm run generate:books.
% canonical-source-hash: fnv1a64:f8cbf64d6d812498
% canonical-lessons: ZH-C124-gezi, ZH-C124-toufa, ZH-C124-shuai, ZH-C124-nankan, ZH-C124-xiang-resemble

\chapter{Tall and Handsome}
\label{ch:zh-tall-and-handsome}
\hlchaptermodality{124}

\begin{chapteropening}
\textbf{By the end of this chapter:} \emph{I can talk about height, hair, good and bad looks, and family likeness.}

Complete the last lesson of chapter 124: I can talk about height, hair, good and bad looks, and family likeness.
\end{chapteropening}

\section[\texorpdfstring{\zh{个子} (gèzi)}{个子 (gèzi)}]{\zh{个子} (gèzi) — height; build (of a person)}

\label{lesson:ZH-C124-gezi}



\begin{quote}
Before the new one: say the Mandarin for warm-hearted, then the Mandarin for to care about.
\end{quote}

\subsection*{You'll want to know: \zh{个子}}
\textbf{\zh{个子}} — \emph{gèzi} — \textquotedblleft{}height; build (of a person)\textquotedblright{}.

\begin{quote}
\emph{wǒ gēge gèzi hěn gāo} — My brother is tall.
\end{quote}

\begin{cousinweb}
\textbf{\zh{个子}} (\emph{gèzi}) is built from \zh{个}, the measure word, and the noun ending \zh{子}, both inherited from Middle Chinese, \emph{kaH} and \emph{tsiX} (Guangyun, in Baxter's transcription). How the pair came to mean a person's height is not settled.
\end{cousinweb}

\begin{grammarlens}[title={What you've built}]
One more word for describing things.
\end{grammarlens}

\subsection*{Your turn}
\begin{itemize}
\raggedright
  \item \emph{Say it:} \emph{gèzi}
  \item \emph{Say it:} \emph{gèzi}, once more
  \item \emph{Say it:} \emph{wǒ gēge gèzi hěn gāo} — My brother is tall.
  \item \emph{Recall:} say the Mandarin for warm-hearted, then the Mandarin for to care about, then say \emph{gèzi} again
\end{itemize}

\begin{tcolorbox}[breakable,colback=teal!4,colframe=teal!35!black,title={Before you move on}]
What does \zh{个子} mean? (\textquotedblleft{}Height; build (of a person)\textquotedblright{}.) Say it once more.
\end{tcolorbox}
\section[tóufa]{tóufa — hair (on the head)}

\label{lesson:ZH-C124-toufa}



\begin{quote}
Before the new one: say the Mandarin for to care about, then the Mandarin for a person's height.
\end{quote}

\subsection*{You'll want to know: tóufa}
\emph{tóufa} — \textquotedblleft{}hair (on the head)\textquotedblright{}.

\begin{quote}
\emph{tā de tóufa hěn cháng} — Her hair is long.
\end{quote}

Its characters come later; for now it is a word for the ear.

\begin{cousinweb}
\emph{tóufa} is built from \emph{tóu} 'head' and \emph{fa}, both inherited from Middle Chinese, \emph{duw} and \emph{pjot} (Guangyun, in Baxter's transcription), which glosses the second 'the hair of the head'. In the word it loses its tone.
\end{cousinweb}

\begin{grammarlens}[title={What you've built}]
One more word for describing things.
\end{grammarlens}

\subsection*{Your turn}
\begin{itemize}
\raggedright
  \item \emph{Say it:} \emph{tóufa}
  \item \emph{Say it:} \emph{tóufa}, once more
  \item \emph{Say it:} \emph{tā de tóufa hěn cháng} — Her hair is long.
  \item \emph{Recall:} say the Mandarin for to care about, then the Mandarin for a person's height, then say \emph{tóufa} again
\end{itemize}

\begin{tcolorbox}[breakable,colback=teal!4,colframe=teal!35!black,title={Before you move on}]
What does tóufa mean? (\textquotedblleft{}Hair (on the head)\textquotedblright{}.) Say it once more.
\end{tcolorbox}
\section[shuài]{shuài — handsome}

\label{lesson:ZH-C124-shuai}



\begin{quote}
Before the new one: say the Mandarin for a person's height, then the Mandarin for hair.
\end{quote}

\subsection*{You'll want to know: shuài}
\emph{shuài} — \textquotedblleft{}handsome\textquotedblright{}.

\begin{quote}
\emph{nǐ gēge hěn shuài} — Your brother is handsome.
\end{quote}

Its characters come later; for now it is a word for the ear.

\begin{cousinweb}
\emph{shuài} is inherited from Middle Chinese \emph{srwijH} (Guangyun, in Baxter's transcription), glossed 'a commander of troops'. The sense handsome is not given there, and how it came is not settled.
\end{cousinweb}

\begin{grammarlens}[title={What you've built}]
One more word for describing things.
\end{grammarlens}

\subsection*{Your turn}
\begin{itemize}
\raggedright
  \item \emph{Say it:} \emph{shuài}
  \item \emph{Say it:} \emph{shuài}, once more
  \item \emph{Say it:} \emph{nǐ gēge hěn shuài} — Your brother is handsome.
  \item \emph{Recall:} say the Mandarin for a person's height, then the Mandarin for hair, then say \emph{shuài} again
\end{itemize}

\begin{tcolorbox}[breakable,colback=teal!4,colframe=teal!35!black,title={Before you move on}]
What does shuài mean? (\textquotedblleft{}Handsome\textquotedblright{}.) Say it once more.
\end{tcolorbox}
\section[nánkàn]{nánkàn — ugly}

\label{lesson:ZH-C124-nankan}



\begin{quote}
Before the new one: say the Mandarin for hair, then the Mandarin for handsome.
\end{quote}

\subsection*{You'll want to know: nánkàn}
\emph{nánkàn} — \textquotedblleft{}ugly\textquotedblright{}.

\begin{quote}
\emph{wǒ de xié hěn nánkàn} — My shoes are ugly.
\end{quote}

Its characters come later; for now it is a word for the ear.

\begin{cousinweb}
\emph{nánkàn} is built from \emph{nán} 'hard' and \zh{看} 'to look at', both inherited from Middle Chinese, \emph{nan} and \emph{khanH} (Guangyun, in Baxter's transcription): hard to look at.
\end{cousinweb}

\begin{grammarlens}[title={What you've built}]
One more word for describing things.
\end{grammarlens}

\subsection*{Your turn}
\begin{itemize}
\raggedright
  \item \emph{Say it:} \emph{nánkàn}
  \item \emph{Say it:} \emph{nánkàn}, once more
  \item \emph{Say it:} \emph{wǒ de xié hěn nánkàn} — My shoes are ugly.
  \item \emph{Recall:} say the Mandarin for hair, then the Mandarin for handsome, then say \emph{nánkàn} again
\end{itemize}

\begin{tcolorbox}[breakable,colback=teal!4,colframe=teal!35!black,title={Before you move on}]
What does nánkàn mean? (\textquotedblleft{}Ugly\textquotedblright{}.) Say it once more.
\end{tcolorbox}
\section[xiàng]{xiàng — to resemble; to take after}

\label{lesson:ZH-C124-xiang-resemble}



\begin{quote}
Before the new one: say the Mandarin for handsome, then the Mandarin for ugly.
\end{quote}

\subsection*{You'll want to know: xiàng}
\emph{xiàng} — \textquotedblleft{}to resemble; to take after\textquotedblright{}.

\begin{quote}
\emph{wǒ xiàng wǒ bàba} — I take after my dad.
\end{quote}

Its characters come later; for now it is a word for the ear.

\begin{cousinweb}
\emph{xiàng} is inherited from Middle Chinese \emph{zjangX} (Guangyun, in Baxter's transcription), glossed 'to be like'.
\end{cousinweb}

\begin{grammarlens}[title={What you've built}]
One more word for describing things.
\end{grammarlens}

\subsection*{Your turn}
\begin{itemize}
\raggedright
  \item \emph{Say it:} \emph{xiàng}
  \item \emph{Say it:} \emph{xiàng}, once more
  \item \emph{Say it:} \emph{wǒ xiàng wǒ bàba} — I take after my dad.
  \item \emph{Recall:} say the Mandarin for handsome, then the Mandarin for ugly, then say \emph{xiàng} again
\end{itemize}

\begin{tcolorbox}[breakable,colback=teal!4,colframe=teal!35!black,title={Before you move on}]
What does xiàng mean? (\textquotedblleft{}To resemble; to take after\textquotedblright{}.) Say it once more.
\end{tcolorbox}
